\documentclass[../../main.tex]{subfiles}% 注意这里的文件路径不能用 ./main.tex ,否则用latexmk编译子文件会报错
\graphicspath{{\subfix{./image/}}} % 指定图片目录,后续可以直接使用图片文件名
% 注意这里的文件路径不能用 ../../image/ ,否则用latexmk编译子文件会报错

% 例如:
% \begin{figure}[H]
% \centering
% \includegraphics[scale=0.3]{图.png}
% \caption{}
% \label{figure:图}
% \end{figure}
% 注意:上述\label{}一定要放在\caption{}之后,否则引用图片序号会只会显示??.

\begin{document}

\section{应用举例}

\begin{definition}
设$S$为正则参数曲面，沿$S$每一点法线截取长度为常数$a$的线段，其端点构成的曲面$S_a$称为$S$的\textbf{平行曲面}。
若两张正则曲面$S,\tilde{S}$之间存在一一对应$\sigma:S\to\tilde{S}$，使得对应点连线为长度常数的公切线，且对应点处两曲面法线夹角为常数，则称该直线族为\textbf{伪球线汇}。
\end{definition}

\begin{example}
设正则曲面$S$的主曲率为$\kappa_1,\kappa_2$，平均曲率为$H$，Gauss曲率为$K$，求其平行曲面$S_a$的主曲率、平均曲率与Gauss曲率。
\end{example}

\begin{solution}
取$S$的二阶标架场$\{\boldsymbol{r};\boldsymbol{e}_1,\boldsymbol{e}_2,\boldsymbol{e}_3\}$，相对分量满足$\omega^3=0,\ \omega_1^3=\kappa_1\omega^1,\ \omega_2^3=\kappa_2\omega^2$，则$S$的基本形式为
\begin{align}
\mathrm{I}=(\omega^1)^2+(\omega^2)^2,\label{eq:paral-I}\\
\mathrm{II}=\kappa_1(\omega^1)^2+\kappa_2(\omega^2)^2.\label{eq:paral-II}
\end{align}
平行曲面$S_a$的向径为$\tilde{\boldsymbol{r}}=\boldsymbol{r}+a\boldsymbol{e}_3$，微分得
\begin{align*}
\mathrm{d}\tilde{\boldsymbol{r}}=\mathrm{d}\boldsymbol{r}+a\mathrm{d}\boldsymbol{e}_3=(\omega^1+a\omega_1^3)\boldsymbol{e}_1+(\omega^2+a\omega_2^3)\boldsymbol{e}_2.
\end{align*}
取$S_a$的一阶标架场$\{\tilde{\boldsymbol{r}};\tilde{\boldsymbol{e}}_1,\tilde{\boldsymbol{e}}_2,\tilde{\boldsymbol{e}}_3\}$，满足$\tilde{\boldsymbol{e}}_1=\boldsymbol{e}_1,\ \tilde{\boldsymbol{e}}_2=\boldsymbol{e}_2,\ \tilde{\boldsymbol{e}}_3=\boldsymbol{e}_3$，则其相对分量为
\begin{align}
\tilde{\omega}^1=(1-a\kappa_1)\omega^1,\quad \tilde{\omega}^2=(1-a\kappa_2)\omega^2,\quad \tilde{\omega}^3=0.\label{eq:paral-omega}
\end{align}
当$a\neq\frac{1}{\kappa_1},\frac{1}{\kappa_2}$时，$S_a$为正则曲面，其第一基本形式为
\begin{align*}
\tilde{\mathrm{I}}=(1-a\kappa_1)^2(\omega^1)^2+(1-a\kappa_2)^2(\omega^2)^2,
\end{align*}
面积元素为
\begin{align}
\mathrm{d}\tilde{\sigma}=\tilde{\omega}^1\wedge\tilde{\omega}^2=(1-2aH+a^2K)\omega^1\wedge\omega^2.\label{eq:paral-area}
\end{align}
由$\mathrm{d}\tilde{\boldsymbol{e}}_3=\mathrm{d}\boldsymbol{e}_3$，得
\begin{align*}
\tilde{\omega}_1^3=-\frac{\kappa_1}{1-a\kappa_1}\tilde{\omega}^1,\quad \tilde{\omega}_2^3=-\frac{\kappa_2}{1-a\kappa_2}\tilde{\omega}^2,
\end{align*}
故$S_a$仍为二阶标架场，其主曲率、平均曲率、Gauss曲率分别为
\begin{align}
\tilde{\kappa}_1=\frac{\kappa_1}{1-a\kappa_1},\quad \tilde{\kappa}_2=\frac{\kappa_2}{1-a\kappa_2},\label{eq:paral-kappa}\\
\tilde{H}=\frac{H-aK}{1-2aH+a^2K},\label{eq:paral-H}\\
\tilde{K}=\frac{K}{1-2aH+a^2K}.\label{eq:paral-K}
\end{align}

\end{solution}

\begin{theorem}[Bäcklund定理]
若欧氏空间$E^3$中两张正则曲面$S,\tilde{S}$由伪球线汇建立一一对应，则$S,\tilde{S}$为具有相同负常Gauss曲率的曲面。
\end{theorem}

\begin{proof}
设伪球线汇中公切线长度为常数$l$，两曲面对应点法线夹角为常数$\tau$。取$S$的一阶标架场$\{\boldsymbol{r};\boldsymbol{e}_1,\boldsymbol{e}_2,\boldsymbol{e}_3\}$，使$\boldsymbol{e}_1$为公切线方向，则$\tilde{\boldsymbol{r}}=\boldsymbol{r}+l\boldsymbol{e}_1$，$\tilde{\boldsymbol{e}}_3=\sin\tau\boldsymbol{e}_2+\cos\tau\boldsymbol{e}_3$。
微分得
\begin{align*}
\mathrm{d}\tilde{\boldsymbol{r}}=\omega^1\boldsymbol{e}_1+(\omega^2+l\omega_1^2)\boldsymbol{e}_2+l\omega_1^3\boldsymbol{e}_3.
\end{align*}
由$\mathrm{d}\tilde{\boldsymbol{r}}\cdot\tilde{\boldsymbol{e}}_3=0$，得
\begin{align*}
\sin\tau(\omega^2+l\omega_1^2)+l\cos\tau\omega_1^3=0,
\end{align*}
即
\begin{align}
\omega_1^2=-\frac{1}{l}\omega^2-\cot\tau\omega_1^3.\label{eq:backlund-omega12}
\end{align}
对式\eqref{eq:backlund-omega12}取外微分，结合结构方程$\mathrm{d}\omega_1^2=-K\omega^1\wedge\omega^2$，整理得
\begin{align*}
-K=\frac{1}{l^2}+\cot^2\tau\cdot K,
\end{align*}
故
\begin{align}
K=-\frac{\sin^2\tau}{l^2}.\label{eq:backlund-K}
\end{align}
由$S,\tilde{S}$地位对称，$\tilde{K}=K$，即两曲面具有相同负常Gauss曲率。

\end{proof}

\begin{proposition}
设$S$为$E^3$中负常Gauss曲率$K=-k^2$的正则曲面，则在其每一点邻域内存在参数系$(u,v)$，使得基本形式为
\begin{align}
\mathrm{I}=\frac{1}{k^2}\left(\cos^2\frac{\varphi}{2}(\mathrm{d}u)^2+\sin^2\frac{\varphi}{2}(\mathrm{d}v)^2\right),\label{eq:sine-I}\\
\mathrm{II}=\frac{1}{k}\cos\frac{\varphi}{2}\sin\frac{\varphi}{2}\left((\mathrm{d}u)^2-(\mathrm{d}v)^2\right),\label{eq:sine-II}
\end{align}
其中$\varphi$为曲面两渐近方向的夹角，满足\textbf{Sine‑Gordon方程}
\begin{align}
\varphi_{uu}-\varphi_{vv}=\sin\varphi.\label{eq:sine-gordon}
\end{align}
反之，若$\varphi$为\eqref{eq:sine-gordon}的解，$k>0$为常数，则存在$E^3$中负常Gauss曲率$K=-k^2$的曲面，以$\varphi$为渐近方向夹角。
\end{proposition}
\begin{proof}
取$S$的正交曲率线网为参数网，基本形式为
\begin{align*}
\mathrm{I}=A^2(\mathrm{d}u)^2+B^2(\mathrm{d}v)^2,\quad \mathrm{II}=\kappa_1A^2(\mathrm{d}u)^2+\kappa_2B^2(\mathrm{d}v)^2,
\end{align*}
取二阶标架场，相对分量为
\begin{align*}
\omega^1=A\mathrm{d}u,\quad \omega^2=B\mathrm{d}v,\quad \omega_1^3=\kappa_1A\mathrm{d}u,\quad \omega_2^3=\kappa_2B\mathrm{d}v.
\end{align*}
设联络形式$\omega_1^2=-\omega_2^1=p\mathrm{d}u+q\mathrm{d}v$，代入结构方程得
\begin{align*}
p=-\frac{A_v}{B},\quad q=\frac{B_u}{A},\quad \omega_1^2=-\frac{A_v}{B}\mathrm{d}u+\frac{B_u}{A}\mathrm{d}v.
\end{align*}
代入Codazzi方程，结合$K=\kappa_1\kappa_2=-k^2$，令$\kappa_1=k\tan\frac{\varphi}{2},\ \kappa_2=-k\cot\frac{\varphi}{2}$，解得
\begin{align*}
A=\cos\frac{\varphi}{2}\cdot a(u),\quad B=\sin\frac{\varphi}{2}\cdot b(v).
\end{align*}
作参数变换$\tilde{u}=k\int a(u)\mathrm{d}u,\ \tilde{v}=k\int b(v)\mathrm{d}v$，仍记为$u,v$，得式\eqref{eq:sine-I},\eqref{eq:sine-II}。
此时联络形式为
\begin{align*}
\omega_1^2=\frac{1}{2}(\varphi_v\mathrm{d}u+\varphi_u\mathrm{d}v),
\end{align*}
代入Gauss方程$\mathrm{d}\omega_1^2=\omega_1^3\wedge\omega_2^3$，即得\eqref{eq:sine-gordon}。
逆命题可由曲面基本定理直接证明。

\end{proof}



\end{document}